% !TeX TS-program = xelatex
% 虚构演示简历 — 仅用于 hellojobs 测试数据
\documentclass{resume-photo}
\usepackage{xeCJK}
\setCJKmainfont{Noto Serif CJK SC}
\usepackage{fontspec}
\usepackage{enumitem}
\setlist[itemize]{itemsep=0pt, topsep=1pt, parsep=0pt, partopsep=0pt}

\ResumeName{陈静}

\begin{document}

\ResumeContacts{
  138-0005-0005,%
  \ResumeUrl{mailto:chenjing@pku.edu.cn}{chenjing@pku.edu.cn},%
  \textnormal{北京大学 | 智能科学与技术 · 硕士 | 2022—2025}%
}

\ResumeTitle

\noindent 大模型应用与 RAG，熟悉 LangChain、向量检索与 LoRA 微调；构建企业内部知识问答，Top-5 命中率离线评测达 81\%。注重可观测性与文档沉淀，习惯在评审阶段拆解风险；参与线上复盘与跨团队联调，英语 CET-6，可阅读官方文档与 RFC（演示数据）。

\section{教育背景}
\ResumeItem
{北京大学}
[\textnormal{智能学院 | 智能科学与技术 | 硕士}]
[2022.09—2025.06]

\begin{itemize}
  \item 导师方向：自然语言处理与知识增强生成。
\end{itemize}


\ResumeItem
{学术交流与在线课程}
[\textnormal{暑校 / MOOC / 厂商认证（演示）}]
[2022—2025]

\begin{itemize}
  \item 完成《深入理解计算机系统》配套实验与 MIT 6.828 / 6.S081 公开课部分章节跟做。
  \item Coursera / edX 分布式系统、云计算相关专项课程结业证书 4 门（虚构编号）。
  \item 通过 AWS SAA / 阿里云 ACP / Redis 开发者等认证中的 1--2 项（简历测试数据）。
  \item 参加 CCF CCSP / 华为 ICT / 百度之星等学科竞赛集训营并完成全部作业。
\end{itemize}



\section{专业技能}
\begin{itemize}
  \item 熟悉 Transformers、PEFT、FAISS/Milvus 向量检索。
  \item 掌握 Prompt 工程、RAG 流水线与评测集构建。
  \item 了解 vLLM 推理服务化与 GPU 显存优化。
  \item 熟悉 Linux 日常运维与 Shell，具备日志检索与 CPU/内存突增初步定位能力。
  \item 掌握 Git / CI 与 Code Review，了解 Docker 与 Kubernetes 基本编排与滚动发布。
  \item 了解 OAuth2 / JWT、接口幂等与 REST 规范，具备单元测试与集成测试习惯（JUnit/pytest 等）。
  \item 能阅读英文技术文档，具备跨团队沟通与需求澄清能力（测试简历）。
\end{itemize}

\section{实习经历}
\ResumeItem
{智谱 AI}
[\textnormal{大模型应用研发实习生}]
[2024.07—2025.01]

\begin{itemize}
  \item \textbf{职责}: 企业客服 Copilot 工具链与评测看板。
  \item \textbf{成果}: 人工抽检一致率从 72\% 提升至 84\%。
\end{itemize}


\ResumeItem
{网易 | 杭州研究院}
[\textnormal{中间件方向实习生}]
[2023.07—2024.01]

\begin{itemize}
  \item \textbf{消息}: 参与 MQ 控制台 Topic 治理与消息轨迹查询性能优化。
  \item \textbf{存储}: 调研并输出对象存储冷热分层方案对比报告 25 页。
  \item \textbf{演练}: 参与双活切换演练，验证 RPO/RTO 指标达标。
  \item \textbf{开源}: 向社区组件提交文档 PR 2 个并被合并。
  \item \textbf{软技能}: 跨 3 团队协调排期，保证里程碑按时交付。
\end{itemize}



\section{项目经历}
\ResumeItem
{多跳问答 RAG}
[\textnormal{课程项目}]
[2024.02—2024.06]

\begin{itemize}
  \item 迭代式检索 + Cross-Encoder 重排，HotpotQA EM 提升 3.8pt。
\end{itemize}


\ResumeItem
{统一配置中心改造}
[\textnormal{Spring Cloud Config → Nacos（演示）}]
[2024.01—2024.04]

\begin{itemize}
  \item \textbf{目标}: 解决配置分散、无灰度与无审计问题，支撑 80+ 微服务。
  \item \textbf{方案}: 双写双读迁移，配置变更走审批流与版本快照。
  \item \textbf{结果}: 配置发布平均耗时从 25min 降至 3min，人为误配下降 60\%。
  \item \textbf{风险}: 设计回滚开关与只读兜底文件，演练 2 次成功。
\end{itemize}



\section{个人荣誉}
\begin{itemize}
  \item ACL 2025 共同作者（在投）
  \item 北京大学 三好学生
  \item 全国计算机等级考试四级（网络工程师）（演示）
  \item 校级「优秀学生干部 / 优秀团员」或技术社团骨干（综合测评前 15\%）
  \item 开源贡献：向知名项目提交被合并的 PR（文档/测试/feature）（演示）
\end{itemize}

\end{document}
